\section{Concurrency}
\label{sec:concurrency}

While taking concurrency into account is necessary
to model most real-world heterogeneous systems,
doing so introduces much complexity, and
models with concurrency are usually
harder to mechanize and reason about.
For example, in sequential programs,
control follows a strict alternation
between the system being modeled and its environment;
this constraint is lost in the presence of concurrency.
The uncertainty associated with scheduling also
introduces a high degree of nondeterminism.

\subsection{Approach}

To manage this complexity,
we will follow a similar approach as we did in \S\ref{sec:effects},
and decouple concurrent layer specifications from
the sequential, deterministic per-thread code that invokes them.
In addition,
using \emph{linearizability} as a correctness criterion
allows us to formulate concurrent specifications in forms that
\begin{itemize}
  \item preserve a strict alternation between
    the system and the environment (atomicity);
  \item factors out scheduling nondeterminism
    into the correctness criterion itself.
\end{itemize}
As such,
the simplest version of our concurrent model
will use a representation for
layer interfaces and implementations
that remains very close to the sequential version,
facilitating the incorporation of refinement conventions and
making sure straightforward specification and verification tasks
can be carried out without undue complexity.

At the same time,
not every concurrent object can be specified
using a deterministic, atomic linearized specification.
To deal with more sophisticated situations,
we will build on the techniques discussed in \S\ref{sec:effects}
to generalize our model by allowing
nondeterministic and non-atomic specifications.

\subsection{Atomic Linearized Specifications}

When multiple threads interact with a concurrent object,
their calls into the object can overlap in unrestricted ways.
The object's behavior in one thread can then be
influenced by any concurrent invocations happening in other threads
while the method is executed.
However,
many concurrent objects behave \emph{as if}
their methods were executed atomically,
at some point between their observable invocation and return.
When that is the case,
a concurrent object can be described
using a specification that remains
essentially sequential and deterministic.
In that case,
the concurrent object implementation
is said to \emph{linearize} to
the sequential atomic specification.

Linearized atomic specifications are
especially useful in our case
as they allow us to avoid much of the complexity associated with
describing concurrent behaviors.
A specification of this kind is essentially
a sequential layer specification as described in \S\ref{sec:effects}
operating over a signature
\newcommand\parr{!}
\[
  \parr E := \{ \alpha . m \mathbin: N \mid
           \alpha \in \Upsilon \wedge
           (m \mathbin: N) \in E \}
  \,,
\]
where each operation is annotated with
a thread identifier $\alpha$
taken from a predefined, fixed set of thread names $\Upsilon$.
In this form,
linearized specifications are much easier to
incorporate in a framework with effects and refinement conventions,
which are inherently sequential in many aspects.

\subsection{Refinement Framework}

With this in mind,
following the strategy established in \S\ref{sec:effects},
a concurrent version of our refinement framework
will be constructed using:
\begin{itemize}
  \item linearized atomic layer specifications as described above;
  \item the sequential (per-thread) implementation strategies
    found in our base framework;
  \item the following notion of layer correctness.
\end{itemize}
We will write
$\Gamma \vdash_\mathbf{R} \sigma : \Delta$
when a per-thread strategy
$\sigma : E \twoheadrightarrow F$,
interacting with an underlay concurrent specification
$\Gamma : M \looparrowright {!E}$,
implements an object which linearizes to
$\Delta : M \looparrowright {!F}$
\emph{up to} the refinement convention $\mathbf{R}$.
Conceptually,
this situation can be depicted as
\begin{equation}
  \Gamma \vdash_\mathbf{R} \sigma : \Delta
  \qquad \qquad
  \begin{tikzcd}
    M \ar[d, equal] \ar[rr, "\Delta"] & &
    !F_1 \ar[d, "\mathbf{R}", leftrightarrow] \\
    M \ar[r, "\Gamma"] &
    !E \ar[r, "!\sigma", twoheadrightarrow] & !F_2
  \end{tikzcd}
  \label{eqn:concref}
\end{equation}
However,
note that in general
the composite $\sigma[\Gamma]$
cannot be defined,
as a \emph{most general linearized atomic specification}
implemented by $\sigma$ may not exist.
In other words,
the property described by (\ref{eqn:concref})
should be seen as a monolithic correctness property
linking the four components $\Delta$, $\Gamma$, $\sigma$ and $\mathbf{R}$.

Nevertheless,
the visualization provided above
is useful to understand how
concurrent refinements compose and
interact with our sequential implementations framework.
In particular,
the following composition principle
can be used to construct
larger layer correctness properties
from elementary ones:
\[
  \begin{prooftree}
    \hypo{\Gamma \vdash_\mathbf{R} \sigma : \Delta}
    \hypo{\Delta \vdash_\mathbf{S} \tau_2 : \Sigma}
    \hypo{\tau_2 \le_{\mathbf{R} \rightarrow \mathbf{T}} \tau_3}
    \infer3{\Gamma \vdash_{\mathbf{S} \mathbin; \mathbf{T}} \tau_3 \circ \sigma : \Sigma}
  \end{prooftree}
  \qquad \qquad
  \begin{tikzcd}
    \bullet \ar[rrr, "\Sigma"] \ar[d, equal] & & &
    !G_1 \ar[d, "\mathbf{S}", leftrightarrow] \\
    \bullet \ar[rr, "\Delta"] \ar[d, equal] & &
    !F_2 \ar[r, "\tau_2"] \ar[d, "\mathbf{R}"', leftrightarrow] &
    !G_2 \ar[d, "\mathbf{T}", leftrightarrow] \\
    \bullet \ar[r, "\Gamma"'] &
    !E_3 \ar[r, "\sigma"'] &
    !F_3 \ar[r, "\tau_3"'] &
    !G_3
  \end{tikzcd}
\]

\subsection{Effects}

Taking an atomic linearized view of concurrent objects
also simplifies incorporating effects into the concurrent framework.
In general,
the way in which effects triggered by different threads may combine
is not obvious.
However,
the sequential ordering induced by linearized atomic specifications
provides an unambiguous answer to this question.

Because of this,
we expect it will be possible to incorporate
effects such as nondeterminism and probabilistic computation
into the concurrent version of the framework
with reasonable effort.
This will enable a variety of important real-world applications,
including in areas such as cryptographic or consensus protocols.

%XXX: could add somthing about how information security
%fits into the picture

\subsection{Progress and Liveness Properties}

While linearized atomic specifications
provide a convenient way to describe the behavior of concurrent objects,
by themselves they can only capture \emph{safety} properties.

When a concurrent object is given a linearized atomic specification,
the object is free to handle
concurrent method invocations by multiple client threads
in an arbitrary order,
introducing a form of demonic nondeterminism.
Unfortunately,
these demonic choices leave open the possibility that
client threads will be starved,
making it hard to provide any guarantees to the client code
based on the linearized atomic specification alone.
To deal with this,
concurrent object specifications are usually augmented
with the choice of a \emph{progress} property
(such as race-freedom, deadlock-freedom, etc)
which constrains the object's behavior in addition to its
linearized atomic specification.
This could be incorporated into our approach
by hard-coding a set number of progress properties
into our definition of layer correctness.

However,
the work of PI Shao and other members of our group show that
a more systematic approach to liveness properties
is possible in the context of linearizability.
We plan to investigate ways in which these techniques
could be incorporated into our paradigm.

\subsection{Shared Memory}

At first blush,
our approach of isolating concurrency
from per-thread language semantics
seems at odd with a treatment of shared-memory concurrency:
in this paradigm,
threads rely only on their local state and
any interactions with other threads
must involve calls to underlay primitives;
whereas accesses to shared memory
would seem to involve a form of cross-thread communication
which by nature must be built into the language semantics.
However,
the combined use of \emph{angelic} nondeterminism
and memory separation constructions
makes it possible to resolve this apparent limitation.

In programs making use of shared-memory concurrency,
accesses to (non-atomic) shared memory regions
must be carefully synchronized between threads to avoid any data races---%
situations where potentially conflicting memory accesses
could lead to undefined behaviors.
As a consequence,
in a race-free execution,
at any point in time shared memory can be decomposed
into individual per-thread memory shares
which assign permissions to individual locations
in a way that is consistent with race-free accesses
across the different threads.
Under this paradigm,
between calls to underlay primitives,
the memory share associated with each thread
can be viewed as a form of private state,
as accesses which are consistent with assigned permissions
cannot interfere with the execution of other threads.

Any actual communication through shared memory
must be carefully synchronized,
for example by acquiring a lock associated with
a particular memory region.
This synchronization will involve
invoking underlay primitives,
which gives us the opportunity to update
the memory share associated with the current thread,
for example by importing exclusive memory permissions
for the region in question.
By incorporating this ``share update'' mechanism
into the specifications of synchronization primitives,
we can bridge the gap between the view of threads as
self-contained sequential computations operating on private state and
the non-atomic shared memory access mechanism
that are required for real-world concurrent programming.

However,
in many low-level languages---%
for example C programs invoking \textsf{pthread} primitives---%
locks and other synchronization mechanisms
are not explicitly linked to the regions of memory they protect,
so that it is impossible for the corresponding primitive specifications
to know in advance how the repartition of memory shares across threads
ought to be updated.
This remaining issue can be addressed using angelic nondeterminism:
under this paradigm, specifications for synchronization primitives
simply enumerate all possible permission updates
which would leave the overall memory unchanged,
in the form of a nondeterministic choice.
In a race-free program,
there will be a way to resolve these choices
such that subsequent memory accesses across all threads will remain valid.
If a proper assignment of permissions if impossible to exhibit,
this signifies the presence of a data race and
will trigger an undefined behavior in the program's semantics.
Since the choice operates on permissions only,
the nondeterminism introduced for distributing shared memory
across thread cannot influence the program's eventual output,
but only whether or not this output is well-defined.

Concretely,
implementing this approach to shared memory
will involve the definition of an appropriate separation algebra
for the CompCert memory model,
as well as incorporating angelic nondeterminism
using the techniques described in the rest of this proposal.
In addition to a treatment of shared memory
which suits our overall paradigm,
we expect this approach to play well with existing separation logic techniques,
and enable us to integrate our framework with
existing tools such as Iris and VST.

%\subsection{Non-atomic Objects}
%
%[Things like interval linearizability, etc.
%All the way up to Arthur-style nonatomic linearization.
%How do we make those fit into a (co)algebraic framework?]

\subsection{Proposed work}

\paragraph{Task 2a: LinCCAL-style correctness with atomic linearized specs}

\paragraph{Task 2b: Liveness properties}

\paragraph{Task 2c: Angelic shared memory}

\paragraph{Task 2d: Non-atomic behaviors}

